\documentclass[a5paper,10pt]{article}

% https://github.com/InDevRus/filippov-solutions
% Нумерация страниц
\pagenumbering{gobble}

% Настройка полей
\usepackage{geometry}
\geometry{left=1cm}
\geometry{right=1cm}
\geometry{top=1cm}
\geometry{bottom=1.5cm}

% Вставка таблицы
\usepackage{array}
\usepackage{tabularx}

% Инструменты выравнивания формул
\usepackage{amsmath}
\usepackage{mathtools}

% Диакритика в математике
\usepackage{amsfonts}

% Вычисления размеров на ходу
\usepackage{calc}

% Удобные записи производных
\usepackage[italic=true]{derivative}

% Настройки сносок
\usepackage{enotez}

% Длинные знаки векторов
\usepackage{anyfontsize}
\usepackage[b]{esvect}

% Вставка картинок
\usepackage{graphicx}

% Вставка красной строки
\usepackage{indentfirst}
\setlength{\parindent}{1em}

% Условия в командах
\usepackage{xifthen}

% Луа-скрипты
\usepackage{luacode}

% Настройка команд отдельно в математическом и текстовом режимах
\usepackage{mathcommand}

% Для повторения LaTeX кода
\usepackage{multido}

% Конфигурация отступов формул
\delimitershortfall-1sp
\usepackage{mleftright}
\mleftright

% Растягивание объектов
\usepackage{relsize}

% Обтекание текстом
\usepackage{wrapfig}
\usepackage{subcaption}
\usepackage{floatflt}

% Поддержка ввода кириллицы
\usepackage[english, russian]{babel}

% Настройка корневой позиции для компилятора
\usepackage{import}
\providecommand*{\ProjectRootPath}{..}

\import{\ProjectRootPath/preambles/macros/}{theorems}

\import{\ProjectRootPath/preambles/fonts}{font_setup}

% Знак градуса
\usepackage{siunitx}

\import{\ProjectRootPath/preambles/macros/}{operators}
\import{\ProjectRootPath/preambles/macros/}{redefinitions}

\import{\ProjectRootPath/preambles/fonts}{letters_setup}
\import{\ProjectRootPath/preambles/fonts/adjustments}{custom_superscripts}

\import{\ProjectRootPath/preambles/custom}{formatting}
\import{\ProjectRootPath/preambles/custom}{frames}
\import{\ProjectRootPath/preambles/custom}{list_setup}

% TODO: перевести команды на современный стиль объявления

\begin{document}

\begin{framed}
    Для угла $\phi = \ang{45}$ имеем $\tg{\beta} = \tg\br{\alpha \pm \phi} =\tg\br{\alpha \pm \ang{45}} = \dfrac {\tg\alpha \pm 1} {1 \mp \tg\alpha }$.
\end{framed}

$\pow{x}{2} + \pow{y}{2} = \pow{a}{2}$.
$2x + 2yy\' = 0$.
Имеем дифференциальное уравнение: $y\' = -\dfrac {x} {y}$.
$\tg\alpha = -\dfrac {x} {y}$.
$\tg\beta$ будет иметь, вообще говоря, два различных значения для заданного $\tg\alpha$.

В случае со знаком <<плюс>>
$\tg \beta = \dfrac {\tg\alpha + 1} {1 - \tg\alpha} = \dfrac {-\dfrac {x} {y} + 1} {1 + \dfrac {x} {y}} = \dfrac {y - x} {x + y} = g\br{x; y\br{x}}$,
а в случае со знаком <<минус>> 
$\tg \beta = \dfrac {\tg\alpha - 1} {1 + \tg\alpha} = \dfrac {-y - x} {y - x} = h\br{x; y\br{x}}$.
Таким образом, имеем два уравнения изогональных траекторий вида $z\'\br{x} = g\br{x; z\br{x}}$ и $z\'\br{x} = h\br{x; z\br{x}}$. Более конкретно, $z\' = \dfrac {z - x} {z + x}$ и $z\' = - \dfrac {z + x} {z - x}$.


\begin{remark}
    Ниже приведены решения соответствующих уравнений в параметрической форме $\tsys{& x\br{t} = C f\br{t} \\& z\br{t} = C t f\br{t}}$ и в полярной форме $\tsys{& x\br{\theta} = C r\br{\theta} \cos\theta \\& z\br{\theta} = C r\br{\theta} \sin\theta}$.

    \centering
    \begin{tabularx}{0.8\textwidth} {
            | >{\centering\arraybackslash}X
            | >{\centering\arraybackslash}X
            | >{\centering\arraybackslash}X |}
        \hline
            Уравнение
            & Функция $f = f\br{t}$
            & Функция $r = r\br{\theta}$ \\
        \hline
            \rule{0mm}{7mm}
            $z\' = \dfrac {z - x} {z + x}$
            & $f\br{t} = \dfrac {\tpow{e}{-\arctg t}} {\sqrt{\pow{t}{2} + 1}}$
            & $r\br{\theta} = \tpow{e}{-\theta}$ \\ [4mm]
        \hline
            \rule{0mm}{7mm}
            $z\' = - \dfrac {z + x} {z - x}$
            & {$f\br{t} = \dfrac {\tpow{e}{\arctg t}} {\sqrt{\pow{t}{2} + 1}}$}
            & $r\br{\theta} = \tpow{e}{\theta}$ \\ [4mm]
        \hline
    \end{tabularx}
\end{remark}

\end{document}

